\chapter{Study and Specification Phase}
\label{chap:etude}

\section{State of the Art and Study of Existing Solutions}
\label{sec:etat-art}

\subsection{Existing Solutions}

Civic issue-reporting platforms are an established category internationally, and reviewing them clarified which patterns to reuse and which gaps to fill for a Tunisia-specific product.

\begin{itemize}[leftmargin=*]
  \item \textbf{FixMyStreet}~\cite{fixmystreet} (UK, mySociety) --- open-source, one of the earliest and most widely deployed civic reporting platforms. Strong on report-to-map visualization and category-based routing; weaker on native mobile experience and not designed for multilingual/RTL usage.
  \item \textbf{311 systems}~\cite{nyc311} (United States --- e.g.\ NYC 311, SeeClickFix) --- large-scale municipal service-request platforms, typically covering more than just infrastructure (they also handle service requests such as noise complaints). They demonstrate the value of a unified tracking code and public status API, both of which Sahali adopts.
  \item \textbf{E-government municipal portals in the MENA region} --- generally administrative-form-oriented (permits, certificates) rather than reactive infrastructure reporting, and rarely offer a native mobile app with offline support.
\end{itemize}

\subsection{Technologies, Methods, and Tools in the Domain}

Across the solutions above, a common technical pattern emerges: a mobile or web front end for citizens, a relational or geospatial database for report storage, a staff-facing dashboard for triage, and a notification layer (SMS/email/push) for status updates. Increasingly, platforms in this space add a machine-learning layer for automatic category classification and duplicate detection --- a pattern Sahali's design anticipates (see the \texttt{ai\_category\_id} and \texttt{is\_duplicate} fields in the data model, Section~\ref{sec:class-diagram}) even though the model itself is future work.

\subsection{Limitations of Existing Solutions}

None of the reviewed solutions fully addresses the constraints this project needed to satisfy:

\begin{itemize}[leftmargin=*]
  \item Limited or no Arabic/RTL support, which is essential for a Tunisian citizen-facing product.
  \item Not designed around Tunisia's municipal structure (multi-tenant by commune, each with its own staff and service-level targets).
  \item Address data in the region is frequently incomplete on OpenStreetMap and available only in one of French or Arabic per point, which existing platforms do not handle explicitly.
  \item Several are either closed-source with licensing costs unsuited to a pilot-stage municipal budget, or open-source but architected for a single large deployment rather than a multi-tenant one.

\end{itemize}

\section{Core Concepts}

This section defines three concepts referenced throughout the report.

\subsection{CI/CD}

\textbf{Continuous Integration (CI)} is the practice of automatically building and testing every code change --- in this project, on every push, via GitHub Actions --- so that regressions are caught within minutes rather than discovered later in production. \textbf{Continuous Deployment/Delivery (CD)} extends this to automating the release process itself: in this project, a signed Android release build is produced automatically from a version tag rather than built by hand on a developer's machine. Together, CI/CD replace manual, error-prone release steps with a repeatable, auditable pipeline (detailed in Section~\ref{sec:cicd}).

\subsection{Cloud Hosting}

\textbf{Cloud hosting} means running infrastructure (servers, databases, storage) on a third-party provider's platform rather than on physical hardware the team owns. Sahali uses three managed cloud services rather than self-hosted infrastructure: \textbf{Render} for the backend API (Docker-based deployment), \textbf{Vercel} for the dashboard (static site hosting with automatic deployment from Git), and \textbf{Supabase} for the production database and file storage (managed PostgreSQL with the PostGIS extension, plus an S3-compatible storage API). This choice removed the need to provision or maintain servers during the internship, at the cost of some platform-specific configuration (e.g.\ passing cryptographic keys as base64-encoded environment variables rather than files).

\subsection{LLM-Assisted Development}

A \textbf{Large Language Model (LLM)} is a machine learning model trained to understand and generate text (and, in coding-assistant contexts, source code) from natural-language prompts. This project was developed with substantial assistance from an LLM-based coding agent (Claude Code), used interactively throughout: for writing and reviewing code, for running and interpreting test suites, for diagnosing build and deployment failures, and for the systematic security audit described in Chapter~\ref{chap:realisation}. This is disclosed here both for transparency and because it directly shaped the iterative, verify-every-step methodology described in Section~\ref{chap:cadre-general} --- an LLM agent can make changes quickly, which raises rather than lowers the importance of verifying each change actually works before moving to the next one.

\section{Specification}

\subsection{Identification of Actors}

\begin{table}[h!]
\centering
\begin{tabularx}{\textwidth}{lX}
\toprule
\textbf{Actor} & \textbf{Role} \\
\midrule
Citizen        & Submits reports (registered or anonymous), tracks their status, receives notifications. \\
Field Agent    & Municipal staff assigned to specific reports; updates status and files a resolution report on completion. \\
Analyst        & Reviews and reclassifies incoming reports, monitors statistics; no administrative rights. \\
Supervisor     & Assigns reports to one or more field agents; manages teams for their municipality. \\
Admin          & Full platform control: municipalities, staff accounts, category/SLA configuration, cross-municipality visibility. \\
\bottomrule
\end{tabularx}
\caption{System actors}
\label{tab:actors}
\end{table}

\subsection{Functional Requirements}

\begin{itemize}[leftmargin=*]
  \item \textbf{Authenticate users} --- registration and login by phone (OTP) or email/password, for both citizens and staff, with role-based access control.
  \item \textbf{Submit a report} --- category, photo, GPS location (with manual correction), optional description; works for both registered and anonymous citizens.
  \item \textbf{Automatic routing} --- match each report to the nearest municipality by geographic distance, with no manual sorting step.
  \item \textbf{Track a report} --- citizens see their own reports' full status history; a public tracking-code lookup works without an account.
  \item \textbf{Triage and assign} --- municipal staff can reclassify a report's category and assign it to one or more field agents.
  \item \textbf{Update report status} --- along the state machine \texttt{submitted → received → under\_review → in\_progress → resolved} (or \texttt{rejected} at any point), enforced server-side.
  \item \textbf{File a resolution report} --- field agents attach up to five photos, an optional video, an optional voice note, and a mandatory written comment when closing a report.
  \item \textbf{Notify} --- push, SMS, and email notifications dispatched on every status change, plus a live in-dashboard event stream for staff.
  \item \textbf{View statistics} --- aggregate counts, resolution-time averages, and per-municipality performance figures on the dashboard.
  \item \textbf{Manage staff and municipalities} --- administrators can create/update/deactivate staff accounts and municipality records.
\end{itemize}

\subsection{Non-Functional Requirements}

\begin{table}[h!]
\centering
\begin{tabularx}{\textwidth}{lX}
\toprule
\textbf{Requirement} & \textbf{Target} \\
\midrule
Multilingual        & Mobile app: Arabic (RTL), French, English. Dashboard: French, Arabic. \\
Responsive           & Dashboard usable on both desktop and tablet widths. \\
Offline capability   & Mobile app queues a submitted report locally (SQLite) and sends it once connectivity returns. \\
Performance          & Report submission end-to-end in under 2 minutes; individual API responses target under 2 seconds. \\
Security             & See the dedicated security requirements below --- treated as a first-class, independently-audited concern rather than folded into general performance targets. \\
Compatibility        & iOS 14+, Android 8.0 (API 26)+; modern evergreen browsers for the dashboard. \\
Maintainability / evolution & Automated test suite and CI pipeline (Chapter~\ref{chap:realisation}) so the codebase can evolve without manual regression testing. \\
\bottomrule
\end{tabularx}
\caption{Non-functional requirements}
\label{tab:nfr}
\end{table}

\textbf{Security requirements} deserve their own enumeration, since they were the subject of a dedicated audit phase (Chapter~\ref{chap:realisation}):
\begin{itemize}[leftmargin=*]
  \item All cryptographic signing keys must be rotatable and never retrievable from source control.
  \item Every endpoint that reads or modifies user data must be authenticated and authorized by role.
  \item Sensitive endpoints (login, OTP, password reset) must be rate-limited independently of the platform-wide default.
  \item Sessions must be revocable server-side (logout must have an immediate effect, not just a client-side token discard).
  \item Concurrent requests against the same security-sensitive resource (a refresh token, a one-time ticket) must not both succeed.

\end{itemize}

\subsection{Constraints}

\begin{itemize}[leftmargin=*]
  \item \textbf{Technical}: the stack was fixed early (Flutter, FastAPI, React/TypeScript, PostgreSQL+PostGIS) based on team familiarity and the need for a geospatial-capable relational database; changing any of these mid-internship was not considered.
  \item \textbf{Temporal}: a summer internship is time-boxed to a few months, which shaped the decision to prioritize a working, secure core platform over building the (larger, higher-risk) AI classification microservice.
  \item \textbf{Material}: no dedicated server budget was available at the start, which drove the choice of managed platforms with usable free/low-cost tiers (Render, Vercel, Supabase) over self-hosted infrastructure.
  \item \textbf{Regulatory}: personal data handling needed to be mindful of Tunisian data protection law~\cite{law632004} (Law 63-2004), which informed decisions such as supporting anonymous report submission.
\end{itemize}

\subsection{General Use Case Diagram}

Figure~\ref{fig:usecase} summarizes the main actor--use-case relationships described above.

\begin{figure}[h!]
\centering
\begin{tikzpicture}[
  actor/.style={minimum height=1.6cm},
  usecase/.style={draw=primary, thick, ellipse, minimum width=3.1cm, minimum height=0.9cm, align=center, font=\small},
  boundary/.style={draw=muted, thick, rounded corners=6pt},
]

  % --- actor drawing macro ---
  \newcommand{\stickfigure}[3]{%
    \begin{scope}[shift={(#1,#2)}]
      \draw[thick] (0,0.85) circle (0.18);
      \draw[thick] (0,0.67) -- (0,0);
      \draw[thick] (-0.28,0.4) -- (0.28,0.4);
      \draw[thick] (0,0) -- (-0.22,-0.45);
      \draw[thick] (0,0) -- (0.22,-0.45);
      \node[font=\small\bfseries, below=0.15cm] at (0,-0.45) {#3};
    \end{scope}
  }

  % Actors (left side)
  \stickfigure{-6.3}{2.6}{Citizen}
  \stickfigure{-6.3}{0.6}{Field Agent}
  \stickfigure{-6.3}{-1.4}{Supervisor}

  % Actors (right side)
  \stickfigure{6.3}{2.0}{Analyst}
  \stickfigure{6.3}{-0.6}{Admin}

  % System boundary
  \node[boundary, minimum width=8.6cm, minimum height=6.4cm] (sys) at (0,0.6) {};
  \node[above, font=\bfseries] at (sys.north) {Sahali Platform};

  % Use cases
  \node[usecase] (uc1) at (-2.6,2.6) {Submit Report};
  \node[usecase] (uc2) at (0,2.6)    {Track Report};
  \node[usecase] (uc3) at (2.6,2.6)  {View Statistics};
  \node[usecase] (uc4) at (-2.6,0.6) {Update Status};
  \node[usecase] (uc5) at (0,0.6)    {File Resolution Report};
  \node[usecase] (uc6) at (2.6,0.6)  {Reclassify Report};
  \node[usecase] (uc7) at (-2.6,-1.4){Assign Report};
  \node[usecase] (uc8) at (0,-1.4)   {Manage Team};
  \node[usecase] (uc9) at (2.6,-1.4) {Manage Staff \& Municipalities};

  % Connections
  \draw (-5.5,2.6) -- (uc1.west);
  \draw (-5.5,2.6) -- (uc2.west);
  \draw (-5.5,0.6) -- (uc4.west);
  \draw (-5.5,0.6) -- (uc5.west);
  \draw (-5.5,-1.4) -- (uc7.west);
  \draw (-5.5,-1.4) -- (uc8.west);

  \draw (5.5,2.0) -- (uc3.east);
  \draw (5.5,2.0) -- (uc6.east);
  \draw (5.5,-0.6) -- (uc9.east);
  \draw (5.5,-0.6) -- (uc7.east);

\end{tikzpicture}
\caption{General use case diagram of the Sahali platform}
\label{fig:usecase}
\end{figure}
